% Фрагмент раздела для модели кривизны через норму градиента.
\section{Граничный минимум и вычисляемая оценка ошибки}
\label{sec:boundary-gradient-model}

Пусть $Q$ есть непустое компактное выпуклое множество, содержащее
более одной точки, функция $f$ принадлежит $C^2$ в окрестности $Q$,
а $x^*\in\arg\min_Q f$. Для всех $z\in Q$ предположим
\begin{equation}\label{eq:boundary-gradient-assumptions}
\begin{aligned}
 \langle\nabla f(z)-\nabla f(x^*),z-x^*\rangle
    &\ge m\|z-x^*\|^2,\qquad m>0,\\
 \|\nabla^2 f(z)\|&\le L_0+L_1\|\nabla f(z)\|,
       \qquad L_0,L_1\ge0.
\end{aligned}
\end{equation}
Пусть $D\ge\operatorname{diam}Q>0$. Полный зазор Франка--Вульфа равен
\[
 G(x)=\max_{v\in Q}\langle\nabla f(x),x-v\rangle.
\]
Первое условие есть сдвинутое условие секущей. Градиент в граничном
минимуме может быть ненулевым. Второе условие ограничивает кривизну
через норму градиента и не требует, чтобы эта норма стремилась к нулю.

\subsection{Каких данных достаточно для оценки ошибки}

\begin{proposition}[Ограничение оценки по норме градиента]
\label{prop:boundary-norm-obstruction}
Зафиксируем $m,D>0$. Пусть сертификат $\mathcal C$ использует только
норму градиента и константы $m,L_0=m,L_1=0,D$, причём для всех задач
с условиями \eqref{eq:boundary-gradient-assumptions} верна оценка
\[
 f(x)-f(x^*)\le\mathcal C(\|\nabla f(x)\|).
\]
Тогда для любого числового аргумента $t\ge0$
\begin{equation}\label{eq:boundary-norm-lower}
 \mathcal C(t)\ge
 \begin{cases}
 t^2/(2m),&0\le t\le mD,\\
 tD-mD^2/2,&t\ge mD.
 \end{cases}
\end{equation}
Поэтому такой сертификат может оставаться сколь угодно большим даже
в точном минимуме функции с постоянной положительной кривизной.
\end{proposition}
\begin{proof}
На фиксированном отрезке $[0,D]$ рассмотрим
$f_a(x)=ax+mx^2/2$, где $a\ge0$. Минимум находится в нуле, вторая
производная равна $m$, а все указанные константы не зависят от $a$.
Если производная в текущей точке равна $t$, то $a=t-mx$ и
\[
 0\le x\le\min\{D,t/m\},\qquad
 f_a(x)-f_a(0)=tx-mx^2/2.
\]
Последнее выражение возрастает на указанном отрезке. Его максимум
равен правой части \eqref{eq:boundary-norm-lower}. В двух случаях он
достигается при $a=0,x=t/m$ и $a=t-mD,x=D$ соответственно.
В точном минимуме функции $f_a$ норма градиента равна $a$, а ошибка
равна нулю, но универсальная оценка $\mathcal C(a)$ всё равно должна
удовлетворять \eqref{eq:boundary-norm-lower}.
\end{proof}

\begin{proposition}[Одного полного зазора недостаточно]
\label{prop:boundary-gap-obstruction}
При условиях \eqref{eq:boundary-gradient-assumptions} с $L_1>0$
значения $G(x),m,L_0,L_1,D$ не задают конечной универсальной верхней
границы для $f(x)-f(x^*)$. Это верно даже на фиксированном квадрате
$[0,1]^2$, когда условие на кривизну выполнено на всей $\mathbb R^2$.
\end{proposition}
\begin{proof}
Положим $Q=[0,1]^2$, $m=1/2$, $L_0=L_1=1$. При $t\ge0$ определим
\[
 f_t(u,v)=\frac{u^2}{2}+\frac t\pi(1-\cos(\pi u))
             +\sqrt{1+v^2}+(\pi t+1)v.
\]
Градиент равен
$(u+t\sin(\pi u),\,\pi t+1+v/\sqrt{1+v^2})$.
Обе его координаты неотрицательны на $Q$, поэтому $x^*=(0,0)$ есть
точка минимума. Кроме того, на $Q$
\[
 \langle\nabla f_t(u,v)-\nabla f_t(0,0),(u,v)\rangle
 =u^2+tu\sin(\pi u)+\frac{v^2}{\sqrt{1+v^2}}
 \ge\tfrac12(u^2+v^2).
\]
На всей $\mathbb R^2$ матрица Гессе равна
$\operatorname{diag}(1+\pi t\cos(\pi u),(1+v^2)^{-3/2})$,
а вторая координата градиента не меньше $\pi t$. Поэтому
\[
 \|\nabla^2 f_t(u,v)\|\le1+\pi t
 \le1+\|\nabla f_t(u,v)\|.
\]
Область и все константы фиксированы. В точке $x=(1,0)$ линейный
оракул выбирает $(0,0)$, поэтому
\[
 G(1,0)=1,\qquad
 f_t(1,0)-f_t(0,0)=\frac12+\frac{2t}{\pi}.
\]
Ошибка неограниченно растёт с $t$, а зазор остаётся неизменным.
\end{proof}

В этом примере допускается отрицательная кривизна вне минимума,
но сдвинутое условие секущей сохраняется. Поэтому к полному зазору
и константам модели нужно добавить другие данные. Один из вариантов
есть норма градиента в текущей точке.

\subsection{Локальная кривизна и явный сертификат}

\begin{theorem}[Сертификат без обнуления нормы градиента]
\label{thm:boundary-local-certificate}
При условиях \eqref{eq:boundary-gradient-assumptions} положим
\begin{align}
 R(x)&=\min\left\{D,\sqrt{G(x)/m},\|\nabla f(x)\|/m\right\},
       \label{eq:boundary-local-radius}\\
 K(x)&=\max\left\{m,
    \bigl(L_0+L_1\|\nabla f(x)\|\bigr)e^{L_1R(x)}\right\},
       \label{eq:boundary-local-curvature}\\
 C(x)&=G(x)+
   \max\left\{0,\frac{K(x)^2-2mK(x)-m^2}{4K(x)}\right\}R(x)^2.
       \label{eq:boundary-local-certificate}
\end{align}
Тогда $f(x)-f(x^*)\le C(x)$ для всех $x\in Q$.
Для любой последовательности в этой фиксированной задаче из
$G(x_k)\to0$ следует $C(x_k)\to0$. Условие
$\|\nabla f(x_k)\|\to0$ не требуется.
\end{theorem}
\begin{proof}
При $x\ne x^*$ положим $r=\|x-x^*\|$ и $e=(x-x^*)/r$.
Из оптимальности на $Q$ следует
$\langle\nabla f(x^*),e\rangle\ge0$. Условие секущей даёт
\[
 \langle\nabla f(x),e\rangle\ge mr,\qquad
 G(x)\ge r\langle\nabla f(x),e\rangle\ge mr^2.
\]
По неравенству Коши--Буняковского также
$mr\le\|\nabla f(x)\|$. Следовательно, $r\le R(x)$.

При движении по отрезку назад от $x$ норма градиента абсолютно
непрерывна и почти всюду удовлетворяет неравенству
\[
 \frac{d}{ds}\|\nabla f(x-se)\|
 \le L_0+L_1\|\nabla f(x-se)\|.
\]
Сравнение со скалярным дифференциальным уравнением даёт
\[
 L_0+L_1\|\nabla f(x-se)\|
 \le\bigl(L_0+L_1\|\nabla f(x)\|\bigr)e^{L_1s},
       \qquad 0\le s\le r.
\]
При $L_1=0$ это утверждение следует сразу из исходной оценки.
Таким образом, $K(x)$ ограничивает норму матрицы Гессе на $[x^*,x]$.

Определим $u(s)=f(x^*+se)-f(x^*)$ на $[0,r]$ и для краткости
положим $K=K(x)$. Полученные оценки дают $|u''|\le K$ и
$u'(0)\le u'(r)-mr$. Поэтому
\[
 u'(s)\le
 \min\{u'(r)-mr+Ks,\ u'(r)+K(r-s)\}.
\]
Поскольку $K\ge m$, две линейные функции пересекаются в точке
$s=(K+m)r/(2K)\in[0,r]$. Интегрирование по двум частям даёт
\begin{equation}\label{eq:boundary-tent-integral}
 u(r)\le r u'(r)
       +\frac{K^2-2mK-m^2}{4K}r^2.
\end{equation}
Здесь $r u'(r)\le G(x)$. Если коэффициент перед $r^2$ неположителен,
отбросим это слагаемое. Иначе применим $r^2\le R(x)^2$.
Получаем \eqref{eq:boundary-local-certificate}. При $x=x^*$
оптимальность на $Q$ даёт $G(x)=0$, и утверждение выполняется сразу.

Непрерывный градиент ограничен на компактном $Q$. Поэтому $K(x_k)$
остаётся ограниченным, а из $G(x_k)\to0$ следует $R(x_k)\to0$.
Правая часть \eqref{eq:boundary-local-certificate} стремится к нулю.
\end{proof}

\paragraph{Применение и ограничения.}
Условие $C(x)\le\varepsilon$ гарантирует нужную точность по функции.
Теорема доказывает сходимость сертификата для последовательности с
исчезающим полным зазором. Сходимость конкретного алгоритма требует
отдельного доказательства. При больших градиентах оценка может быть
грубой. Если дополнительно известна проверенная граница
$\|\nabla^2f(z)\|\le K_Q$ на $Q$, величину $K(x)$ можно заменить на
$\min\{K(x),K_Q\}$. Приведённый контрпример исключает универсальную
оценку ошибки только через $m,L_0,L_1,D$ и зазор. Количественные
оценки работы, получаемые через такую границу ошибки, требуют
дополнительных данных, например границы градиента или кривизны
и начальной ошибки. Сам контрпример не даёт нижней оценки числа шагов.
Точность полученного явного сертификата как верхней границы для
всего класса с кривизной через градиент не установлена.
